\documentclass[10pt,a4paper]{amsart}
\usepackage[margin=19mm]{geometry}
\usepackage{amsmath,amssymb,mathtools}
\usepackage[scr=rsfs,cal=euler]{mathalfa}
\usepackage{xfrac}
\usepackage[unicode,hidelinks,bookmarks=false]{hyperref}
\newcommand{\ie}{i.e.\ }
\newcommand{\cf}{cf.\ }
\newcommand{\resp}{respectively\ }
\newcommand{\Subsection}{\subsection}
\providecommand{\nobreakdash}{\nobreak\hbox{-}}
\setlength{\emergencystretch}{2em}
\multlinegap0pt
\usepackage{mathtools}
\usepackage{dsfont}
\usepackage{bm}
\newtheorem{assumption}{Assumption}
\newtheorem{lemma}{Lemma}
\newtheorem{theorem}{Theorem}
\title{\LARGE \bf Distributed Vector Assembly via Selector-Anchored Diffusion with Local Innovation}
\author{Diana Vieira Fernandes, Carlos Santos Silva}
\date{Source: arXiv:2609.10201v2 (CC BY 4.0)}
\begin{document}
\maketitle

\noindent\emph{Source and licence.} \LARGE \bf Distributed Vector Assembly via Selector-Anchored Diffusion with Local Innovation. Version arXiv:2609.10201v2 (CC BY 4.0), distributed by arXiv under the Creative Commons Attribution 4.0 International licence, http://creativecommons.org/licenses/by/4.0/, as recorded in the arXiv licence metadata for this exact version and shown on the abstract page https://arxiv.org/abs/2609.10201v2. Full licence notice: \texttt{licenses/arxiv-2609.10201-CC-BY-4.0.txt}.\par\smallskip

\noindent\textit{Editorial scope.} Counted unit: the article's lemma lem:l2\_gain (source0.tex lines 1258--1287). The article's complete original proof of it is delivered with it (source0.tex lines 1289--1367, one \texttt{proof} environment of 79 lines). No mathematics is written, repaired or re-derived: the statement and the proof are the article's own lines, in the article's own notation, and no retained line is substituted or re-flowed.  Retained uncounted as the source-local context the proof needs: lines 147--150; lines 213--219; lines 309--320; lines 343--353.  Nothing was omitted from any retained span, so the package carries no omission record.  No retained line contains a citation, so the wrapper carries no bibliography and no bibliography entry.  No retained line contains a figure environment, an image-inclusion command or drawing code, so no image is delivered and the package declares no asset dependency.  Editorial formatting only, in addition: an AMS \texttt{amsart} preamble that restates the article's own declarations and macros used by the retained lines, namely the environments \texttt{assumption}, \texttt{lemma}, \texttt{theorem} on separate counters, together with the standard packages those lines need.  The retained environments print in this document as Assumption 1, Lemma 1, Theorem 1 in order of appearance, whereas the article prints the same items with the numbering of its own full document.  The complete native source of the article as distributed in the arXiv e-print for this exact version is bundled unchanged as \texttt{sources/209883/source0.tex}.
\begin{assumption}[Communication graph]
\label{ass:graph}
The graph \(\mathcal G\) is fixed, connected, and undirected.
\end{assumption}

\begin{equation}\label{eq:selector_partition}
\mathbf{E}_i^\top \mathbf{E}_i = I_{d_i},
\quad
\mathbf{E}_i^\top \mathbf{E}_j = \mathbf{0}\quad \forall i\neq j,
\quad
\sum_{i=1}^N \mathbf{E}_i \mathbf{E}_i^\top = I_d.
\end{equation}

\begin{lemma}[Global Contraction]\label{lem:contraction}
Under Assumption~\ref{ass:graph}, the Laplacian \(L=L^\top\) has eigenvalues \(0=\lambda_1(L)<\lambda_2(L)\le\cdots\le\lambda_N(L)\).
Let $\mathbf{P}=\mathrm{blkdiag}(\mathbf{P}_1,\dots,\mathbf{P}_N)$ with $\mathbf{P}_i=\mathbf{E}_i\mathbf{E}_i^\top$. By the selector-partition identity \eqref{eq:selector_partition}, $\sum_{i=1}^N \mathbf{P}_i = I_d$. Define $W := I_N - \alpha L$ and
\begin{align}
\mathbf{A}_{\mathrm{sys}} \;:=\; W\otimes I_d \;-\; \mathbf{P}.
\end{align}
If $0<\alpha<1/\lambda_N(L)$, then $\|\mathbf{A}_{\mathrm{sys}}\|_2<1$. Consequently, the mapping is a contraction in the Euclidean norm:
\begin{align}
\|\mathbf{A}_{\mathrm{sys}}\mathbf{x}\|_2 \le \sigma \|\mathbf{x}\|_2,\qquad \forall \mathbf{x}\in\mathbb{R}^{Nd},
\end{align}
with $\sigma:=\|\mathbf{A}_{\mathrm{sys}}\|_2\in[0,1)$.
\end{lemma}

\begin{theorem}[Finite-$T$ contraction and outer-step ISS]\label{thm:iss_2ts}
Under the hypotheses of Lemma~\ref{lem:contraction}, let $\sigma:=\|\mathbf{A}_{\mathrm{sys}}\|_2\in[0,1)$. Then, for every integer $T\ge1$, the inner loop contracts as
\begin{equation}\label{eq:inner_contr_T}
\|\mathbf{e}^{k,T}\|_2 \le \sigma^{T}\,\|\mathbf{e}^{k,0}\|_2.
\end{equation}
Moreover, the outer-step propagation satisfies
\begin{equation}\label{eq:outer_iss}
\|\mathbf{e}^{k+1,0}\|_2
\le \sigma^{T}\,\|\mathbf{e}^{k,0}\|_2 + \sqrt{N}\,\|\Delta u^{k}\|_2.
\end{equation}
\end{theorem}

\begin{lemma}[Estimator $\ell_2$-gain]\label{lem:l2_gain}
Let \(a:=\sigma^T\in[0,1)\). Under the hypotheses of Theorem~\ref{thm:iss_2ts},
\begin{align}
\|\mathbf e^{k+1,0}\|_2
&\le
a\|\mathbf e^{k,0}\|_2+\sqrt N\|\Delta u^k\|_2, \notag\\
\|\mathbf e^{k,T}\|_2
&\le
a\|\mathbf e^{k,0}\|_2 .
\end{align}
Then, for any \(\nu>0\) and any \(K\ge 1\),
\begin{align}
\sum_{k=0}^{K-1}\|\mathbf e^{k,T}\|_2^2
&\le
(1+\nu)\frac{a^2}{1-a^2}\|\mathbf e^{0,0}\|_2^2 \notag\\
&\quad +
\Gamma_e(T,\nu)
\sum_{k=0}^{K-1}\|\Delta u^k\|_2^2,
\end{align}
where
\begin{align}
\Gamma_e(T,\nu)
:=
(1+\nu^{-1})
\frac{N a^2}{(1-a)^2}
=
(1+\nu^{-1})
\frac{N\sigma^{2T}}{(1-\sigma^T)^2}.
\end{align}
\end{lemma}

\begin{proof}
Define
\begin{equation}
x_k:=\|\mathbf e^{k,0}\|_2,
\qquad
d_k:=\sqrt N\,\|\Delta u^k\|_2 .
\end{equation}
The estimator recursion gives
\begin{equation}
x_{k+1}\le a x_k+d_k .
\end{equation}
Iterating this inequality yields, for every \(k\ge 0\),
\begin{equation}
x_k
\le
a^k x_0
+
\sum_{j=0}^{k-1}a^{k-1-j}d_j,
\end{equation}
where the sum is understood to be zero when \(k=0\). Since \(\|\mathbf e^{k,T}\|_2\le a x_k\), it follows that
\begin{equation}
\|\mathbf e^{k,T}\|_2
\le
a^{k+1}x_0
+
\sum_{j=0}^{k-1}a^{k-j}d_j .
\end{equation}
For any \(\nu>0\), Young's inequality gives
\begin{align}
\|\mathbf e^{k,T}\|_2^2
\le{}&
(1+\nu)a^{2k+2}x_0^2 \notag\\
&+
(1+\nu^{-1})
\left(
\sum_{j=0}^{k-1}a^{k-j}d_j
\right)^2 .
\end{align}
Summing the first term over \(k=0,\dots,K-1\) gives
\begin{equation}
\sum_{k=0}^{K-1}a^{2k+2}x_0^2
\le
\frac{a^2}{1-a^2}x_0^2 .
\end{equation}

For the second term, consider the causal convolution kernel \(h_\ell:=a^\ell\), \(\ell\ge1\). Its \(\ell_1\)-norm is
\begin{equation}
\|h\|_{\ell_1}
=
\sum_{\ell=1}^{\infty}a^\ell
=
\frac{a}{1-a}.
\end{equation}
Young's convolution inequality therefore gives
\begin{align}
\sum_{k=0}^{K-1}
\left(
\sum_{j=0}^{k-1}a^{k-j}d_j
\right)^2
&\le
\left(\frac{a}{1-a}\right)^2
\sum_{k=0}^{K-1}d_k^2 \notag\\
&=
\frac{Na^2}{(1-a)^2}
\sum_{k=0}^{K-1}\|\Delta u^k\|_2^2 .
\end{align}
Combining the preceding bounds yields
\begin{align}
\sum_{k=0}^{K-1}\|\mathbf e^{k,T}\|_2^2
\le{}&
(1+\nu)\frac{a^2}{1-a^2}
\|\mathbf e^{0,0}\|_2^2 \notag\\
&+
(1+\nu^{-1})
\frac{Na^2}{(1-a)^2}
\sum_{k=0}^{K-1}\|\Delta u^k\|_2^2,
\end{align}
which is the claimed estimate.
\end{proof}

\end{document}
